%% ma: L11-B16-0053
%% lop: 11
%% chuong: 5
%% bai: 16
%% dang: 11.16.2
%% loai: TN4
%% muc: VDC
%% dap_an: "B"
%% nguon: "tai-lieu-hoc-tap-mon-toan-lop-11-theo-chuong-trinh-giao-duc-pho-thong-2018.pdf (Lớp Toán Thầy Duy, Phạm Lê Duy, THPT Chu Văn An), Chương 3 Bài 2, Phần 1, Câu 41 (tr.372)"
%% kien_thuc: "Lớp 11 Bài 16 (giới hạn của hàm số)"
%% trang_thai: chinh-thuc
%% ghi_chu: "Đã duyệt 03/11/2026."
\begin{ex}
Cho $f(x)$ là đa thức thỏa mãn $\lim\limits_{x\to2}\dfrac{f(x)-20}{x-2}=10$. Tính $T=\lim\limits_{x\to2}\dfrac{\sqrt[3]{6f(x)+5}-5}{x^2+x-6}$.
\choice{$T=\dfrac{12}{25}$}{\True $T=\dfrac4{25}$}{$T=\dfrac4{15}$}{$T=\dfrac6{25}$}
\loigiai{\sloppy Từ giả thiết, $f(2)=20$ (đa thức liên tục) và $\dfrac{f(x)-20}{x-2}\to10$. Đặt $A=\sqrt[3]{6f(x)+5}$: $A-5=\dfrac{6f(x)+5-125}{A^2+5A+25}=\dfrac{6\left(f(x)-20\right)}{A^2+5A+25}$. Mà $x^2+x-6=(x-2)(x+3)$. Do đó \[T=\lim\dfrac{6}{(x+3)\left(A^2+5A+25\right)}\cdot\dfrac{f(x)-20}{x-2}=\dfrac{6\cdot10}{5\cdot75}=\dfrac4{25}.\]}
\end{ex}
